\hypertarget{twenty-steerable-individually-switchable-tweezers-at-10644-and-8946-nm}{%
\chapter{Twenty steerable, individually switchable tweezers at 1064.4
and 894.6
nm}\label{twenty_tweezers-twenty-steerable-individually-switchable-tweezers-at-10644-and-8946-nm}}

Design study for Objective 3 (proposal ANRF/ARG/2026/009937, draft v97).
Code: \texttt{picsim/opa/star.py}, \texttt{scripts/m8\_twenty.py};
results \texttt{results/m8\_*.json}. All simulated values are SIM
outputs, valid in the regime simulated, and must be reviewed before they
enter the proposal.

\begin{quote}
Follow-up: this design moves all lit tweezers together. For independent
motion of each tweezer, as with AOD tones, see
\texttt{docs/independent\_steering.md} (sliding-window focal-plane OPA).
The switched-port design below remains the recommendation for the 894.6
nm addressing beams.
\end{quote}

\hypertarget{1-result}{%
\section{Result}\label{twenty_tweezers-1-result}}

Twenty tweezers that are steerable and individually switchable can be
delivered on the 400 nm SiN platform by a multi-port star-coupler OPA
with 128 emitters, one chip per wavelength. Each tweezer has its own
input port and its own MZI switch; the ports illuminate a common 128-arm
phased array through a free-propagation region (FPR); a phase ramp on
the arms moves all lit tweezers together. On simulated designs, with the
20 ports placed on a row of 4 um sites:

{\small
\begin{longtable}[]{@{}>{\raggedright\arraybackslash}p{0.355\dimexpr(\linewidth-6\tabcolsep)\relax}>{\raggedright\arraybackslash}p{0.372\dimexpr(\linewidth-6\tabcolsep)\relax}>{\raggedright\arraybackslash}p{0.273\dimexpr(\linewidth-6\tabcolsep)\relax}@{}}
\toprule\noalign{}
& 1064.4 nm (transport tweezers) & 894.6 nm (addressing beams) \\
\midrule\noalign{}
\endhead
\bottomrule\noalign{}
\endlastfoot
spot FWHM along the row & 1.52 um & 1.27 um \\
intensity at the 4 um neighbour / peak & 1.0e-6 & 3.5e-6 (neighbour
error 8.5e-6) \\
peak spread, all 20 lit, random port phases & \textless= 0.73\% &
\textless= 0.89\% \\
stray at unlit sites (10 of 20 lit) & \textless= 1.2e-5 & \textless=
2.1e-5 \\
static position error of a spot & \textless= 6 nm & \textless= 5 nm \\
switch per port & one trimmed thermo-optic MZI & two cascaded AlN
MZIs \\
extinction needed & \textgreater= 31 dB (1 uK residual at 1.2 mK) &
\textgreater= 50 dB (idle error 2.5e-5) \\
fibre power & 33-43 mW per tweezer at 1.2 mK; 0.74 W for 20 & \textless=
0.14 mW per site; 2.4 mW for 20 \\
\end{longtable}
}

The price is that steering is collective: all lit tweezers move by the
same displacement. Which sites carry tweezers is set by which ports are
lit. This is the primitive an AOD provides for parallel rearrangement
(select a subset of a row, move it), but at one optical frequency, so
two tweezers brought together do not beat. Fully independent steering of
20 tweezers through one aperture requires phase-only synthesis on a
shared OPA, and that architecture is a poor on/off engine (Sec. 4).

\hypertarget{2-architecture}{%
\section{Architecture}\label{twenty_tweezers-2-architecture}}

One chip per wavelength, both on the shared 400 nm LPCVD SiN process:

(i) edge-coupled input (two fibres at 1064.4 nm, to keep each facet
below 20 MW/cm\^{}2); (ii) a 1x32 tunable-splitter tree (five stages of
thermo-optic MZIs with 3 um 2x2 MMIs), which sends power only to lit
ports, so the power per tweezer does not depend on how many are lit;
(iii) 20 port switches, thermo-optic at 1064.4 nm (tweezer on/off ramps
of \textgreater= 100 us), two cascaded AlN MZIs at 894.6 nm (gate
pulses, ZVD-shaped drive, \textless= 100 ns); (iv) 20 input tapers, 1.4
um (1064.4 nm) or 1.0 um (894.6 nm) wide, spaced 1.89 um and 1.64 um at
the FPR; (v) an FPR of radius 223 um (1064.4 nm) or 197 um (894.6 nm);
(vi) 128 arms, 1.7 um wide at the FPR edge (0.3 um gaps), each with a
phase shifter, routed to (vii) 128 end-fire emitters at 1.5 um pitch;
(viii) the relay of Sec. 6 of \texttt{docs/findings.md}: cylindrical
collimation, Fourier-plane slit, anamorphic telescope, meta-objective
(NA 0.5).

The two chips are combined before the meta-objective by a dichroic.
Their port geometries map to the same 4 um sites; the residual
registration is set by the common phase ramp of each chip.

Each port illuminates all 128 arms with the same envelope, the far field
of its taper mode, which acts as a Gaussian-like apodisation of the
array. Port j imposes a linear phase k\_s d\_j p\_a / R per arm (k\_s,
the slab wavenumber; d\_j, port offset; p\_a, arm pitch), which steers
its beam to sin(theta\_j) = n\_s d\_j p\_a / (R p\_out). The port pitch
that places beams 4 um apart at the atoms follows from the relay scale x
= sin(theta) N p\_out / (2 NA): a 4 um step is sin(theta) = 0.0208 for N
= 128 and p\_out = 1.5 um.

\hypertarget{3-design-choices-and-trade-offs}{%
\section{Design choices and
trade-offs}\label{twenty_tweezers-3-design-choices-and-trade-offs}}

\textbf{Apodisation.} The arm-power envelope sets the spot tails, and
the tails set two things: the light that a lit tweezer leaves on its
neighbours, and the interference between lit ports, whose relative
phases drift with the router temperature. With the envelope truncated at
1/e\^{}2 at the outer arms (as for a standard star coupler), the
neighbour intensity is 1.2e-3 and, with all 20 lit and random port
phases, peak intensities vary by up to 19\%. Narrowing the envelope
removes both at the cost of a wider spot:

{\small
\begin{longtable}[]{@{}>{\raggedright\arraybackslash}p{0.337\dimexpr(\linewidth-8\tabcolsep)\relax}>{\raggedright\arraybackslash}p{0.185\dimexpr(\linewidth-8\tabcolsep)\relax}>{\raggedright\arraybackslash}p{0.272\dimexpr(\linewidth-8\tabcolsep)\relax}>{\raggedright\arraybackslash}p{0.207\dimexpr(\linewidth-8\tabcolsep)\relax}@{}}
\toprule\noalign{}
envelope at outer arm (nominal) & FWHM at 1064.4 nm & neighbour / peak &
peak spread, 20 lit \\
\midrule\noalign{}
\endhead
\bottomrule\noalign{}
\endlastfoot
e\^{}-2 & 1.10 um & 1.2e-3 & 19\% \\
e\^{}-4 & 1.24 um & 4.2e-4 & 11\% \\
e\^{}-6 & 1.38 um & 9.6e-5 & 5.8\% \\
e\^{}-8 & 1.52 um & 1.3e-5 & 2.0\% \\
\end{longtable}
}

(Input taper 0.6 um.) A narrower envelope needs a smaller FPR radius,
which brings the ports closer together (1.0 um at e\^{}-8). Widening the
input taper restores the radius: with 1.4 um tapers the port pitch is
1.89 um, and the taper mode, being closer to Gaussian than the 0.6 um
guide mode, also lowers the tails further. Two 1.4 um tapers 0.49 um
apart have a coupling length of 1.33 mm at 1064.4 nm, so the
\textasciitilde20 um over which they converge on the FPR transfers well
below -20 dB. The final designs were chosen from a scan of taper width
and envelope (\texttt{results/m8\_taper\_scan.json}): 1.4 um tapers at
1064.4 nm, 1.0 um tapers at 894.6 nm, both at a nominal e\^{}-6
envelope. Because the wide-taper modes are not Gaussian, the realised
envelope at the outer arm is 2.7e-5 (1064.4 nm) and 7.2e-6 (894.6 nm),
and a nominally narrower envelope (e\^{}-8) puts the outer arms on the
side lobes of the port far field (realised 1e-3 to 2.5e-3, worse tails).
The designs are therefore specified by FPR radius and taper width, not
by the nominal edge level.

We keep the cleaner design at 1064.4 nm although a wider envelope
(e\^{}-5) gives 9\% more depth per watt, because its 3\% peak spread and
4e-5 neighbour light would load the static atoms and make tweezer depths
depend on drifting port phases.

\textbf{Emitter pitch.} 1.5 um, the smallest pitch at which a 50 um
straight end-fire section keeps emitter-to-emitter transfer below -20 dB
(\texttt{m3\_pitch\_efficiency}). It raises the slit efficiency to 0.65
from 0.52 at 2 um.

\textbf{Channel count.} With N = 96 the steering half-range is 51 um
(1064.4 nm) or 43 um (894.6 nm), so a row of 20 sites (+/-38 um) leaves
only 5 to 13 um for collective moves; N = 128 gives +/-68 um and +/-57
um.

\hypertarget{4-why-not-phase-synthesis-on-one-opa}{%
\section{Why not phase synthesis on one
OPA}\label{twenty_tweezers-4-why-not-phase-synthesis-on-one-opa}}

On a shared 128-channel OPA, 20 tweezers can be computed with
independent positions (\texttt{m3\_multibeam}), but a tweezer is
switched off only by re-synthesis. Removing one of 20 (warm-started
refinement with that site among the dark points) leaves 1.9 to 3.3\% of
the tweezer intensity at its site (-15 to -17 dB), changes the 19 kept
tweezers by up to 12\%, and along a linear phase path between the two
patterns the kept tweezers deviate by up to 19\%. Stray light at
untargeted sites stays at about 2\%. Hence it cannot gate 894.6 nm light
at 50 dB, and it perturbs the other atoms whenever one tweezer is
switched. It remains the route to fully independent steering at 1064.4
nm if that is ever required.

\hypertarget{5-steering}{%
\section{Steering}\label{twenty_tweezers-5-steering}}

All lit tweezers move by the displacement set by a linear phase ramp
across the 128 arms. With all 20 lit (peak change quoted at 1064.4 nm;
at 894.6 nm it is within 0.6\% of these values):

{\small
\begin{longtable}[]{@{}lll@{}}
\toprule\noalign{}
collective shift & position error (static, max) & peak change (min /
max) \\
\midrule\noalign{}
\endhead
\bottomrule\noalign{}
\endlastfoot
4 um & 25 nm (1064.4), 5 nm (894.6) & -2.5\% / +2.3\% \\
8 um & 25 nm, 5 nm & -5.1\% / +4.3\% \\
16 um & 25 nm, 6 nm & -10.7\% / +7.8\% \\
24 um & 26 nm (1064.4 only) & -16.7\% / +10.3\% \\
\end{longtable}
}

The position error is a static offset of the outer spots, caused by the
element-factor slope across them; it changes by less than 2 nm over a 24
um sweep and is calibrated out. The peak change follows the emitter
element factor and can be pre-compensated through the port powers. At
894.6 nm the outer spots leave the +/-57 um field beyond a 19 um shift;
subsets of the row can move farther.

Because the collective move is one linear ramp, its fast part needs a
single waveform: the 128 arm shifters take fixed gains n = 0 ... 127
from one analog ramp, with slow per-arm offsets for calibration. A
continuous analog ramp has no staircase, which removes the update-rate
heating of \texttt{m3\_update\_rate} (40 quanta per 100 us move at 100
kHz and 1.2 mK). The shifters must be fast and low-power: PZT on SiN
(0.4 MHz, V\_pi about 5.6 V at 1064.4 nm by wavelength scaling of
{[}MONT{]}) rather than thermo-optic (34 us rise).

\hypertarget{6-switching}{%
\section{Switching}\label{twenty_tweezers-6-switching}}

At 1064.4 nm a tweezer is ramped on or off adiabatically over
\textgreater= 100 us, so a thermo-optic MZI (34 us rise) suffices;
\textgreater= 31 dB keeps the residual below 1 uK at 1.2 mK, which a
trimmed single MZI exceeds at phase errors up to 10 mrad (5th percentile
38 dB, \texttt{m4\_extinction}; computed at 894.6 nm, to be repeated at
1064.4 nm). At 894.6 nm the switch gates Raman pulses. The
proposal\textquotesingle s 50 dB gives an idle-qubit error of 2.5e-5;
the switch leak equals the 3.5e-6 optical tail of a lit neighbour at
54.6 dB. Two cascaded AlN MZIs give a 5th percentile of 66 to 88 dB, so
the switch is not the limit; ZVD-shaped drive reaches -50 dB within 49
ns for a 2\% error in the actuator resonance (\texttt{m4\_dynamics},
with an assumed resonance).

\hypertarget{7-power}{%
\section{Power}\label{twenty_tweezers-7-power}}

At 1064.4 nm, per tweezer, from the fibre: facet -0.37 dB, router -0.28
dB, switch -0.06 dB, propagation (1 cm, 0.3 dB/cm assumed) -0.30 dB,
FPR, arms, emitters, slit and pupil 0.40 to 0.52 of the port power
depending on the port, meta-objective 0.7 (assumed): -6.3 to -7.4 dB to
the atom plane. For a 1.2 mK depth (1160 a.u., Gaussian 1.04 um along
the cross axis) this needs 33 to 43 mW at the fibre per tweezer, 0.74 W
for 20, against 1.4 W in the proposal. At 894.6 nm, 20 sites at the
intensity of 20 uW in a 1.04 um Gaussian need 2.4 mW in total.

\hypertarget{8-implementation-load}{%
\section{Implementation load}\label{twenty_tweezers-8-implementation-load}}

{\small
\begin{longtable}[]{@{}>{\raggedright\arraybackslash}p{0.348\dimexpr(\linewidth-6\tabcolsep)\relax}>{\raggedright\arraybackslash}p{0.339\dimexpr(\linewidth-6\tabcolsep)\relax}>{\raggedright\arraybackslash}p{0.312\dimexpr(\linewidth-6\tabcolsep)\relax}@{}}
\toprule\noalign{}
item & 1064.4 nm chip & 894.6 nm chip \\
\midrule\noalign{}
\endhead
\bottomrule\noalign{}
\endlastfoot
router MZIs (thermo-optic) & 31 & 31 \\
port switches & 20 thermo-optic & 40 AlN (two per port, high voltage) \\
arm phase shifters & 128 PZT (one fast ramp + slow offsets) & 128 (slow;
registration only) \\
FPR radius & 223 um & 197 um \\
electrical channels & about 180 & about 200 \\
\end{longtable}
}

About 380 channels across the two chips exceed the 32-64-channel OPA and
16-channel router of the proposal, and the 32-channel high-voltage
driver of Sec. 8 covers only the 894.6 nm switches if one AlN MZI per
port is replaced by a trimmed thermo-optic stage (a trimmed single MZI
reaches 50 dB in 89\% of devices at 3 mrad). Thermo-optic power, at half
of P\_pi on average, is about 42 mW per heater without trenches
(\texttt{m5\_heaters}), i.e. about 2.1 W for 51 heaters on the 1064.4 nm
chip; trenches with undercut bring this to about 0.1 W.

\hypertarget{9-what-this-does-not-provide-and-limits-of-the-model}{%
\section{What this does not provide, and limits of the
model}\label{twenty_tweezers-9-what-this-does-not-provide-and-limits-of-the-model}}

It does not move different tweezers by different amounts in one step; a
sorting sequence is built from parallel moves, as with AODs. It steers
along one row; the second axis needs a second chip, a row-select
element, or an OPA-AOD hybrid (\texttt{docs/findings.md}, Sec. 6). The
FPR model is 2D and paraxial (flat arcs, no Rowland aberration, isolated
arm modes), and the arm phases from fabrication are assumed calibrated.
The relay is scalar and aberration-free; a meta-objective with a Strehl
of 0.8 adds a halo that is not included in the crosstalk figures.

\hypertarget{10-suggested-wording-for-the-proposal}{%
\section{Suggested wording for the
proposal}\label{twenty_tweezers-10-suggested-wording-for-the-proposal}}

Hypothesis H3b: "An OPA fed through a multi-port star coupler forms up
to 20 tweezers, each switched on and off by its own integrated MZI, and
moves all lit tweezers together by setting a single phase ramp across
its 128 emitters. All tweezers share one optical frequency, hence ... ."
Table 1, OPA row: "simultaneous, individually switchable tweezers
\textgreater= 20 on one row; collective move \textgreater= +/-16 um with
all 20 lit; neighbour intensity \textless= 1e-4 at 894.6 nm".

\hypertarget{verification-notes}{%
\section{Verification notes}\label{twenty_tweezers-verification-notes}}

Sources read in session: as listed in \texttt{docs/findings.md}
({[}MONT{]} arXiv:2601.15695 for the PZT figures). No new literature
values were introduced in this study; all new numbers are SIM or
DERIVED.

{\footnotesize
\begin{longtable}[]{@{}>{\raggedright\arraybackslash}p{0.302\dimexpr(\linewidth-8\tabcolsep)\relax}>{\raggedright\arraybackslash}p{0.141\dimexpr(\linewidth-8\tabcolsep)\relax}>{\raggedright\arraybackslash}p{0.151\dimexpr(\linewidth-8\tabcolsep)\relax}>{\raggedright\arraybackslash}p{0.407\dimexpr(\linewidth-8\tabcolsep)\relax}@{}}
\toprule\noalign{}
value & location & type & pointer \\
\midrule\noalign{}
\endhead
\bottomrule\noalign{}
\endlastfoot
FWHM, neighbour, spread, stray, position error, delivery & Sec. 1, 3, 5
& SIM & results/m8\_final.json (review before submission) \\
apodisation table & Sec. 3 & SIM & results/m8\_apod.json (review) \\
taper-width scan, realised envelope 2.7e-5 / 7.2e-6 & Sec. 3 & SIM &
results/m8\_taper\_scan.json, m8\_final.json (review) \\
coupling length 1.33 mm (0.49 um gap, 1.4 um tapers) & Sec. 3 & SIM &
femwell, run in session; noted in m8\_final.json \\
re-synthesis residual 1.9-3.3\%, 12\%, 19\% & Sec. 4 & SIM &
results/m8\_resynth.json (review) \\
collective-shift table & Sec. 5 & SIM & results/m8\_final.json, steer
(review) \\
33-43 mW, 0.74 W, 2.4 mW & Sec. 1, 7 & SIM/DERIVED &
results/m8\_final.json, power (review) \\
31 dB, 50 dB, 54.6 dB & Sec. 6 & DERIVED & results/m8\_final.json,
switching \\
emitter pitch 1.5 um; arm pitch 2.0 um; arm width 1.7 um; propagation
0.3 dB/cm; meta-objective 0.7 & Sec. 2, 7 & CHOICE & stated
assumptions \\
1.2 mK, 1160 a.u., 1.04 um, 4 um sites, NA 0.5, 20 uW per site &
throughout & USER & proposal draft v97 \\
\end{longtable}
}

Open placeholders: {[}VALUE?{]} PZT shifter footprint and hysteresis for
128 arms; {[}VALUE?{]} AlN resonance and Q at 894.6 nm; {[}VALUE?{]}
measured arm phase errors after fabrication (calibration range).
